Models with the same stable-fate predictions can differ in whether a response
persists after a signal is withdrawn. We compare the published Calzone
death-receptor model with its documented CASP3-to-CASP8 feedback-deletion
variant, using one fixed interface observing NFkB, CASP3 and mitochondrial
permeability transition. Both variants reach survival, apoptosis and necrosis
under sustained TNF. An automatically extracted shared state follows
\DRDepth{} asynchronous updates and the visible history CASP3 activation.
From that state, sustained-stimulation continuations have the same exact
observable language. Allowing TNF withdrawal breaks correspondence: the
feedback-positive continuation preserves CASP3 and has only an apoptotic
terminal fate, whereas the feedback-negative continuation permits CASP3
deactivation and has only the model's naive terminal signature. Exact
simulation certificates and an independent model checker verify this
conditioned comparison. A verified 12-action counterexample also refutes
global sustained trace equivalence; the reverse trace inclusion remains
undecided. Thus branching alone does not separate globally trace-equivalent
biological models in this case. A secondary
DNA-damage example illustrates correspondence changes across three declared
interfaces. The study provides a reproducible relational reconstruction of a
published commitment phenomenon, not a newly discovered mechanism or wet-lab
validation.
